\begin{theorem}[Wild quadratic common correction]
\label{mod:cases-wildquadratic-commoncorrection}
Let $K/F$ be wildly ramified quadratic with residue degree one, let $t$ be
its lower break, and put
\[
 T=t+1=v_K(\mathfrak D_{K/F}),\qquad m=m_F(\chi_F).
\]
All data below are relative to one of the simultaneous representative
packages supplied by the high- or low-parameter theorem.  No representative
is declared canonical.

In the low range $m\leq T$, retain the supplied
$\epsilon _1,\alpha _1,\beta _1$ and their exact norms
$\epsilon,\alpha,\beta$.  The low adapter uses that same selected pair and
sets
\[
 u=\frac{\epsilon _1\beta _1}{\alpha _1},
 \qquad n=\frac{\epsilon\beta}{\alpha}.
\]
In the high range $T\leq m$, retain the entire supplied high package,
including
\[
 \delta_h\in U_F^t,
 \qquad \beta=\delta_hN_{K/F}(\beta _1).
\]
Writing $\alpha=N_{K/F}(\alpha _1)$ and $c=\alpha\delta_h$, the high adapter
sets
\[
 u=\frac{\beta _1}{\alpha _1},
 \qquad n=\frac{\beta}{\alpha\delta_h}=\frac\beta c.
\]
Lean proves the literal identities $cn=\beta$ and
\[
 N_{K/F}(\text{the selected high correction term})
   =\beta(\alpha\delta_h).
\]
Thus the essential HighQuadratic unit $\delta_h$ remains visible in the
normalization and its norm identity.

Both adapters return the same choice-relative record type $C=(u,n)$.  The
record stores the exact facts
\[
 [K:F]=2,\qquad f(K/F)=1,\qquad
 v_K(\mathfrak D_{K/F})=T,
\]
\[
 v_K(u)=v_F(n)=T-m,\qquad N_{K/F}(u)=n,\qquad 1+n\ne0.
\]
For the low boundary $m=T$, the last fact is proved using the actual nonzero
stationary quotient class of the selected $j=1$ twist.  Lean uses the
already selected field representative and does not manufacture or declare a
canonical lift.
The same actual twist noncancellation, now combined with residue
characteristic two, also excludes $n=1$ for both adapters.  Consequently
$1+u$ and $1+u^\vee$ are nonzero, so Lean exposes $z_0$ and $z_1$ as
elements of $F^\times$ for multiplicative-character evaluation.  The generic
unit constructors require an explicit opposite-noncancellation proof, and
each stationary adapter supplies precisely that proof.

For one such record define in the fields
\[
 s=\operatorname{Tr}_{K/F}(u),\qquad
 u^\vee=-\frac n u,
\]
\[
 z_0=\frac{N_{K/F}(1+u^\vee)}{1+n},\qquad
 z_1=\frac{N_{K/F}(1+u)}{1+n},
 \qquad x=z_0-1,\quad y=z_1-1,
\]
and $X=s+nx+y$.  Lean proves the exact identities
\[
 N(u^\vee)=n,\qquad \operatorname{Tr}(u^\vee)=-s,
\]
\[
 N(1+u)=1+s+n,\qquad N(1+u^\vee)=1-s+n,
\]
and, with the manuscript's signs and denominator direction,
\begin{equation}\label{eq:lean-wild-quadratic-common-X}
 x=-\frac{s}{1+n},\qquad
 y=\frac{s}{1+n},\qquad
 X=s+nx+y=\frac{2s}{1+n}.
\end{equation}
It also proves the literal field equality
\[
 \operatorname{Tr}_{K/F}(u^{-1})=\frac{s}{n}.
\]
For the four selected stationary parameters, Lean additionally preserves
the exact ratio and additive identities
\[
 n-u=-u(1+u^\vee),\qquad
 N(n-u)=n(1+n)z_0,\qquad N(1+u)=(1+n)z_1,
\]
\[
 \operatorname{Tr}(n-u)+1-n-(n+1)=-s.
\]

Suppose $m=2d+1>T$ and put $r=m-T$.  The selected $u^{-1}$ has exact
order $r$, and the ramification API is used at precisely the adjacent
lattice depths
\[
 \mathfrak p_K^r/\mathfrak p_K^{r+1}
   \xrightarrow{\ \sim\ }
 \mathfrak p_F^d/\mathfrak p_F^{d+1}.
\]
For $c_\chi=-\operatorname{Tr}_{K/F}(u^{-1})$, Lean proves
$v_F(c_\chi)=d$ and the graded-quotient equality
\[
 \operatorname{grTr}_{r,d}([-u^{-1}])=[c_\chi].
\]
It does not turn this class equality into a new equality of field
representatives.  The associated formula-facing equality is literal:
\[
 \frac{x}{c_\chi}=\frac{n}{1+n}.
\]

At the odd boundary $m=T=2d_\tau+1$, the already selected $u$ is bundled as
an integral unit.  For every explicitly supplied
$\widetilde\rho\in\mathcal O_F$ whose residue maps to the residue of this
$u$, Lean proves
\[
 \operatorname{Tr}_{K/F}(u)-2\widetilde\rho
   \in\mathfrak p_F^{d_\tau+1}
\]
and hence, only in the displayed quotient,
\[
 \operatorname{grTr}_{0,d_\tau}([u])=[2\widetilde\rho]
 \quad\text{in }\mathfrak p_F^{d_\tau}/\mathfrak p_F^{d_\tau+1}.
\]
The residue lift remains an explicit argument.

The ambiguity API has exactly the strength of the parameter APIs.  For two
low packages $P,Q$ with the same fixed $\epsilon_1$, it retains the target
and norm congruences for $\alpha$ at depth $r$ and for $\beta$ at depth $d$,
the source ratios in $U_K^r$ and $U_K^d$, and only the derived statement
\[
 \frac{u(P)}{u(Q)}\in U_K^{\min(r,d)}.
\]
For two high packages it retains precisely the proved quotient equalities
for $\beta$ modulo $\mathfrak p_F^d$, for
$c=\alpha\delta_h$ in
$\mathfrak p_F^{m-T}/\mathfrak p_F^{m-T+\lfloor T/2\rfloor}$, for the twist
representative modulo $\mathfrak p_F^d$, for the upstairs representative
modulo $\mathfrak p_K^{d_K}$, and for its norm-product class modulo
$\mathfrak p_F^d$.  It proves no literal equality between the two choices
of $u,n,s,x,y,$ or $X$, and it does not identify arbitrary high and low
representatives at the boundary.

The principal theorem bundles the exact norm expansions and all three
equalities in \eqref{eq:lean-wild-quadratic-common-X} for one common record.
The surrounding exported declarations provide the two explicit adapters,
the retained high unit correction, the two graded-trace specializations,
and the exact ambiguity interfaces.  No coefficient-table, error-formula,
boundary-cancellation, or First Main Lemma theorem is assumed.

\LeanFile{LanglandsFirstMainLemma/Cases/WildQuadratic/CommonCorrection.lean}
\lean{LanglandsFirstMainLemma.WildQuadraticCommonCorrectionData}
\lean{LanglandsFirstMainLemma.lowWildQuadraticCommonCorrection}
\lean{LanglandsFirstMainLemma.highWildQuadraticCommonCorrection}
\lean{LanglandsFirstMainLemma.highWildQuadraticCommonCorrection_n_with_delta}
\lean{LanglandsFirstMainLemma.highWildQuadraticCommonCorrection_delta_factorization}
\lean{LanglandsFirstMainLemma.highWildQuadraticCommonCorrection_exact_norm}
\lean{LanglandsFirstMainLemma.lowWildQuadraticCommonCorrection_oppositeNoncancellation}
\lean{LanglandsFirstMainLemma.highWildQuadraticCommonCorrection_oppositeNoncancellation}
\lean{LanglandsFirstMainLemma.LowWildQuadraticCommonCorrectionAmbiguity}
\lean{LanglandsFirstMainLemma.lowWildQuadraticCommonCorrection_ambiguity}
\lean{LanglandsFirstMainLemma.HighWildQuadraticCommonCorrectionAmbiguity}
\lean{LanglandsFirstMainLemma.highWildQuadraticCommonCorrection_ambiguity}
\lean{LanglandsFirstMainLemma.WildQuadraticCommonCorrectionData.norm_one_add_u}
\lean{LanglandsFirstMainLemma.WildQuadraticCommonCorrectionData.norm_one_add_dual}
\lean{LanglandsFirstMainLemma.WildQuadraticCommonCorrectionData.norm_n_sub_u_eq}
\lean{LanglandsFirstMainLemma.WildQuadraticCommonCorrectionData.norm_one_add_u_eq_denominator_mul_z1}
\lean{LanglandsFirstMainLemma.WildQuadraticCommonCorrectionData.signed_trace_sum_eq}
\lean{LanglandsFirstMainLemma.WildQuadraticCommonCorrectionData.z0Unit}
\lean{LanglandsFirstMainLemma.WildQuadraticCommonCorrectionData.z1Unit}
\lean{LanglandsFirstMainLemma.WildQuadraticCommonCorrectionData.X_eq}
\lean{LanglandsFirstMainLemma.WildQuadraticCommonCorrectionData.trace_inv_eq}
\lean{LanglandsFirstMainLemma.wildQuadraticHighOdd_gradedTrace_class}
\lean{LanglandsFirstMainLemma.wildQuadraticHighOdd_cChi_order}
\lean{LanglandsFirstMainLemma.wildQuadraticHighOdd_x_div_cChi_eq}
\lean{LanglandsFirstMainLemma.wildQuadraticBoundary_trace_congr}
\lean{LanglandsFirstMainLemma.wildQuadraticBoundary_gradedTrace_class}
\lean{LanglandsFirstMainLemma.wildQuadratic_commonCorrection}
\leanok
\SourceText{Lemma lem:quad-X and Lemma lem:quadratic-graded-trace.}
\uses{mod:parameters-highquadratic,mod:parameters-low,mod:ramification-normfiltration}
\FormalizationContract
The high and low constructions consume their supplied simultaneous
existential packages; they do not construct canonical field
representatives.  The high package retains $\delta_h$.  Norm, trace, and
correction equations are literal field equalities only after one package is
fixed.  The strict-odd and boundary trace conclusions are equalities in the
stated adjacent lattice quotients, and comparisons between packages use
only the quotient or filtration precision certified upstream.
\end{theorem}
